\chapter{Proponowane tematy projektów}

\begin{summarybox}
Ten rozdział nie zastępuje oficjalnej listy zadań laboratoryjnych. Jest katalogiem
pomysłów i wariantów projektów, które pomagają dobrać temat do poziomu studenta.
Każdy temat powinien zostać przekształcony w konkretny raport: model, równania,
kod, wykresy, kontrola wyniku i interpretacja.
\end{summarybox}

\section{Jak wybierać temat?}

Dobry temat powinien być wystarczająco prosty, aby dało się go ukończyć, ale
wystarczająco bogaty, aby wymagał interpretacji. Nie należy wybierać tematu tylko
dlatego, że generuje ładny wykres. Projekt ma pokazać rozumienie mechaniki
kwantowej i umiejętność pracy obliczeniowej.

Przy wyborze tematu warto odpowiedzieć na pytania:
\begin{itemize}
    \item jaki model fizyczny będzie badany,
    \item jakie równanie zostanie rozwiązane,
    \item co będzie liczone albo wizualizowane,
    \item jak zostanie sprawdzona poprawność wyniku,
    \item jaka będzie główna interpretacja fizyczna.
\end{itemize}

\section{Tematy łatwe}

\begin{exercisebox}
\textbf{Nieskończona studnia potencjału.}
Zbadać energie, funkcje własne i gęstości prawdopodobieństwa dla kilku pierwszych
stanów. Pokazać zależność energii od szerokości studni. Kontrola: normalizacja i
zgodność z analitycznym wzorem na \(E_n\).
\end{exercisebox}

\begin{exercisebox}
\textbf{Normalizacja funkcji falowych.}
Porównać normalizację funkcji prostokątnej, sinusoidalnej i Gaussowskiej. Pokazać,
jak zmienia się gęstość prawdopodobieństwa po zmianie parametrów. Kontrola:
całka z \(|\psi|^2\) równa jedności.
\end{exercisebox}

\begin{exercisebox}
\textbf{Oscylator harmoniczny: pierwsze stany.}
Narysować funkcje własne i gęstości prawdopodobieństwa dla kilku pierwszych
stanów oscylatora. Porównać energię kwantową z intuicją klasyczną. Kontrola:
orthogonalność i normalizacja.
\end{exercisebox}

\section{Tematy średnie}

\begin{exercisebox}
\textbf{Skończona studnia potencjału.}
Znaleźć numerycznie energie stanów związanych dla studni o zadanej głębokości i
szerokości. Porównać wynik z nieskończoną studnią. Kontrola: liczba stanów
związanych i zachowanie funkcji falowej poza studnią.
\end{exercisebox}

\begin{exercisebox}
\textbf{Tunelowanie przez barierę prostokątną.}
Zbadać zależność współczynnika transmisji od energii, wysokości i szerokości
bariery. Pokazać przypadki \(E<V_0\) i \(E>V_0\). Kontrola: \(R+T\approx1\).
\end{exercisebox}

\begin{exercisebox}
\textbf{Paczka Gaussa padająca na barierę.}
Rozwiązać zależne od czasu równanie Schrödingera dla paczki falowej. Pokazać
odbicie i transmisję w czasie. Kontrola: norma całkowita oraz oddzielenie paczki
odbitej i transmitowanej.
\end{exercisebox}

\begin{exercisebox}
\textbf{Atom wodoru: orbitale i radialna gęstość prawdopodobieństwa.}
Narysować wybrane orbitale albo przekroje \(|\psi|^2\). Porównać \(|R(r)|^2\) z
\(r^2|R(r)|^2\). Kontrola: normalizacja radialna i poprawne liczby kwantowe.
\end{exercisebox}

\section{Tematy trudniejsze}

\begin{exercisebox}
\textbf{Model Kroniga--Penneya.}
Zbadać powstawanie pasm i przerw energetycznych w potencjale periodycznym.
Pokazać dozwolone i zabronione zakresy energii. Kontrola: analiza warunku
Blocha i porównanie z pojedynczą barierą.
\end{exercisebox}

\begin{exercisebox}
\textbf{Poziomy Landaua.}
Zbadać cząstkę w jednorodnym polu magnetycznym i pokazać redukcję do oscylatora
harmonicznego. Obliczyć odstępy poziomów dla wybranych wartości pola. Kontrola:
jednostki i zależność od \(B\).
\end{exercisebox}

\begin{exercisebox}
\textbf{Grafen w modelu ciasnego wiązania.}
Zbudować funkcję \(f(\mathbf{k})\), narysować strukturę pasmową i wskazać punkty
Diraca. Kontrola: sprawdzenie, gdzie \(f(\mathbf{k})=0\), oraz analiza liniowej
dyspersji w pobliżu tych punktów.
\end{exercisebox}

\begin{exercisebox}
\textbf{Sferyczna kropka kwantowa.}
Wyznaczyć energie w sferycznej studni potencjału przez zera funkcji Bessela
sferycznych. Narysować wybrane przekroje gęstości prawdopodobieństwa. Kontrola:
warunek brzegowy i degeneracja \(2l+1\).
\end{exercisebox}

\section{Tematy rozszerzone}

Tematy rozszerzone mogą obejmować skończoną sferyczną studnię potencjału,
absorpcję na brzegach domeny w TDSE, split-operator FFT, zaburzenia oscylatora,
model podwójnej studni, ewolucję kubitu pod działaniem Hamiltonianu dwupoziomowego
albo porównanie obliczeń Mathematica/Python dla tego samego problemu.

Takie tematy są dobre dla studentów, którzy dobrze opanowali część podstawową.
Powinny jednak nadal zawierać kontrolę wyniku, a nie tylko bardziej złożony kod.

\section{Minimalne wymagania dla każdego projektu}

Każdy projekt, niezależnie od poziomu trudności, powinien zawierać:
\begin{itemize}
    \item opis modelu,
    \item równanie albo Hamiltonian,
    \item wartości parametrów i jednostki,
    \item kod obliczeniowy,
    \item co najmniej dwa sensowne wykresy,
    \item kontrolę wyniku,
    \item interpretację fizyczną,
    \item krótkie wnioski.
\end{itemize}

\begin{summarybox}
Najlepszy projekt nie musi być najbardziej efektowny. Najlepszy projekt jest taki,
w którym wiadomo, jaki model badamy, co zostało policzone, jak sprawdzono wynik i
jaki jest sens fizyczny otrzymanych wykresów.
\end{summarybox}
